\documentclass[11pt]{article}
\usepackage[margin=1in]{geometry}
\usepackage{amsmath,amssymb,amsthm,hyperref}
\newtheorem{lemma}{Lemma}
\newcommand{\OPT}{\operatorname{OPT}}
\setlength{\emergencystretch}{3em}
\title{A fixed-partition projection lemma and its remaining access gap}
\author{Research note, 2 October 2026}
\date{}
\begin{document}
\maketitle

\noindent\textbf{Status.}
This is a hand lemma about one fixed primary preferred-path partition.
It does not supply a cheap comparator for the evolving multi-splay forest.
Its purpose is to isolate the rotation part of that possible comparison
from the access and split/join parts. In particular, it makes no claim
that an arbitrary projection of a global search path is cheap.

Let $P$ be a fixed full leaf-oriented reference tree. Its internal
boundaries are the actual keys; its leaves are the external gaps.
Fix one preferred child at each internal node. A nonempty component
$C$ is a maximal preferred path, continued to its terminal external
leaf $z_C$. Include empty terminal components at nonpreferred external
children when discussing the component tree. Write $S_C$ for the
reference subtree beginning at $C$'s first node; for an empty component
this is its one external leaf.

The exit blocks of $C$ are its terminal leaf and the reference subtrees
at its nonpreferred children. Each latter subtree starts one child
component $D$ in the component tree. Define
\[
 W_C=\{z_C\}\cup\{z_D:D\text{ is a child component of }C\}.
\]
There is exactly one representative in each exit block. If $C$ has
$k$ internal nodes, then $W_C$ has $k+1$ leaves. Every external leaf
belongs to the representative sets of at most two components: its own
terminal component and that component's parent.

For an arbitrary global routing tree $T$ on the external leaves, let
$T\upharpoonright W_C$ be its induced tree: prune the other leaves and
suppress unary nodes. Relabel the $k$ resulting internal boundaries
by the actual primary separators of $C$ in increasing key order.
This gives a BST on exactly the component's actual key set.

\begin{lemma}[Two components meet three consecutive intervals]
Let $A,B,D$ be nonempty disjoint leaf intervals in this order. At most
two nonempty components $C$ satisfy
\[
 W_C\cap A\ne\varnothing,\qquad
 W_C\cap B\ne\varnothing,\qquad
 W_C\cap D\ne\varnothing.
\]
The intervals need not cover the whole reference universe.
\end{lemma}

\begin{proof}
For any such component, its reference domain contains a leaf of $A$
and a leaf of $D$, hence contains every leaf of $B$. Domains of distinct
such components are therefore nested, and their components form a
chain in the component tree.

Suppose $C$ is above another such component $E$. Let $F$ be the
immediate child component of $C$ whose domain contains $S_E$. All
representatives of $C$ other than $z_F$ lie outside $S_F$, hence
outside $S_E$. Because $B\subseteq S_E$ and $W_C$ meets $B$, it follows
that $z_F\in B$. If $E$ were strictly below $F$, then $E$ would lie
inside a nonpreferred child subtree off the preferred path from $F$
to $z_F$. Its domain would exclude $z_F$, contradicting
$z_F\in B\subseteq S_E$. Thus $E=F$ and its terminal leaf lies in $B$.

No still lower component can have a domain containing all of $B$:
each of its first nonpreferred edges below $E$ excludes $z_E\in B$.
Consequently the chain has at most two members.
\end{proof}

\begin{lemma}[Constant total rotation congestion]
For the fixed preferred partition above, one rotation of $T$ changes
the collection $\{T\upharpoonright W_C\}_C$ by at most two rotations
in total. Relabeling their boundaries by actual component keys does
not change this conclusion.
\end{lemma}

\begin{proof}
Write a routing rotation as $((A,B),D)\leftrightarrow(A,(B,D))$,
where the three hanging subtrees have disjoint consecutive leaf
intervals. If a representative set misses one of those three intervals,
pruning and unary suppression give the same induced tree on both
sides. If it meets all three, the two induced trees differ by one
rotation. The previous lemma bounds the latter sets by two.

The selected leaf sets are fixed during this operation, so the ordered
component separator labels are fixed as well.
\end{proof}

Thus a fixed partition admits a projection of a comparator's entire
rotation trace with at most twice as many rotations. No potential or
component-lifetime initialization charge is involved in this statement.
Initial projected trees may be installed in total linear cost because
their disjoint actual key sets contain $n$ keys in total.
For root or constant-depth global rotations the projected rotations
also have constant rooted pointer cost, since pruning shortens ancestor
paths. The rotation count alone does not assert this for an arbitrary
source trace with a finger deep inside its current tree.

\paragraph{Why this does not yet prove a comparator theorem.}
An original query $q$ may lie in one exit block $E$ while the selected
representative $z_E$ lies elsewhere in that block. The global search
for $q$ can take a different branch from the search for $z_E$ at a
global separator inside $E$. Relabeling the induced tree's boundaries
does not identify these two searches.

For example, take a nonextreme exit block whose neighboring component
keys are $\ell<r$. Put the global root at a requested key $q\in E$,
and choose $z_E>q$. Arrange the representative leaves to the right of
$q$ so that their induced routing tree is a long left spine ending at
$z_E$. After projection, the root boundary is $\ell$, while $r$ is
deep on that left spine. The global search for $q$ costs one, but the
component search into its $E$-gap, or exposure of its upper guard $r$,
can have linear cost in the component size. A mirrored example gives
the other direction. Such a component and gap can occur on a monotone
primary preferred path, with $E$ one of its nonpreferred exits.

This is a single-state obstruction to a pointwise access comparison,
not a recurrent lower bound. A global amortized argument could still
pay the exposure. Moreover, changing primary preferences changes the
sets $W_C$. The rotation lemma does not pay for replacing representatives,
separating an old tail, or attaching an incoming tail under the exact
common-root split/join checkpoints. Those costs must be proved before
one can infer a bound such as
$C_{\mathrm{forest}}(W_{\mathrm{MS}})=O(\OPT+n)$.

\paragraph{A separate priority-based observation.}
There is another way to avoid arbitrary depth ties in a projection.
Give a global BST distinct heap-ordered priorities. If a child $y$ is
rotated above its parent $x$, decrease only the priority of $y$ to a
value strictly between the priorities of $x$ and its former parent
(or below the root priority if $x$ was the root). All other priorities
may remain unchanged, and the new global tree is still its Cartesian
tree. Thus Cartesian restrictions to a fixed key partition can change
only in the component containing $y$.

This last observation bounds the number of affected components, not
the work in that component: $y$ may pass many projection ancestors that
were in the opposite subtree of $x$. It is therefore not a substitute
for the two-rotation leaf-projection lemma or for the missing amortized
access argument.

\section*{An endpoint refinement for fixed-partition accesses}

The access obstruction above has a useful partial repair. For every
exit block $E$ of a component $C$, select both its minimum and maximum
external leaves, selecting a singleton only once. Call the resulting
set $V_C$ and prune $T$ to $V_C$.

\begin{lemma}[Endpoint congestion]
A global routing rotation changes at most one of the induced trees
$T\upharpoonright V_C$.
\end{lemma}

\begin{proof}
Again a changing induced tree must have a selected leaf in each of
the three consecutive rotation intervals $A,B,D$. Two changing
components would have nested reference domains, each containing all
of $B$. If $C$ is above $F$, its only selected points inside the child
domain leading to $F$ are the two endpoints of that domain. Because
this child domain contains a leaf of $A$ and a leaf of $D$, both of
its endpoints lie outside $B$. The other representatives of $C$ are
outside that whole domain. This contradicts $V_C\cap B\ne\varnothing$.
\end{proof}

The internal boundaries of $T\upharpoonright V_C$ alternate between
real component separators and at most one additional, or ghost,
separator inside each nonsingleton exit block $E$. The ghost node for
$E$ is exactly the global lowest common ancestor of its two selected
endpoints. Its global key lies inside $E$.

Convert this induced tree to an alphabetic routing tree on the exit
blocks as follows. A real separator supplies the usual two-way split.
A ghost node for $E$ supplies a three-way ordered split into lower
blocks, $E$, and upper blocks; refine this into two binary comparisons
at $E$'s adjacent real component boundaries, omitting an absent side.
An external selected leaf carries the query block only when that block
was a singleton; otherwise its ghost node carries it. Suppress empty
and unary refinements. Each original node contributes at most two
physical routing nodes. Every block occurs once, so each real component
separator occurs exactly once.

This conversion has the same rooted-hull property as the no-copy
unfolding: a fixed abstract frontier has fixed ordered block intervals,
so a change of $k$ abstract hull nodes changes at most $2k$ physical
nodes with exactly the same hanging components. In particular, a
projected rotation costs $O(1)$ rotations after conversion. Together
with endpoint congestion, root or constant-depth global rotations have
$O(1)$ total rooted projection cost.

If an actual query key $q$ lies in a nonsingleton exit block $E$, its
global search visits the global lowest common ancestor of $E$'s two
endpoint leaves. Therefore the converted component search to its
$E$-gap has length at most twice the global search length. If $q$ is
itself a component separator, its projected physical image is reached
no later than the corresponding twice-length bound. This remedies the
single-representative access obstruction for a fixed component.

One can also sum the searches over the nested components encountered
on a reference search for $q$. The projected route for component $C$
uses global search nodes in $S_C\setminus E$, followed by the first
global node inside $E$, where $E$ is the next encountered component's
domain. The former sets are disjoint corridors. The latter contributes
one additional node per component. Thus the total projected search
length is $O(\operatorname{depth}_T(q)+b)$ for $b$ encountered components.

\paragraph{The remaining opposite guard.}
This endpoint refinement still does not prove the evolving-forest
comparison. At a left-to-right primary switch, access to the incoming
right gap exposes $y$ and its upper ancestor guard $r$. The lower
ancestor guard $\ell$ lies beyond all keys of the old left tail and
need not be exposed by this access. The exact common forest word
also prescribes its split at $\ell$.

Collapsing the whole old tail to one new exit block would expose the
other gap's two boundaries $\ell,y$, but changes the selected endpoint
set and the projection itself. The fixed-partition congestion lemma
does not pay for that coarsening, nor for transporting the resulting
projection states through the six prescribed split/join checkpoints.
Proving an amortized bound for these changes remains necessary; a
pointwise assertion that the incoming-gap search already supplies all
three standard multi-splay guards would be false.

\end{document}
